
\documentclass[11pt,a4paper,oneside]{article}

\usepackage[a4paper,margin=2.2cm]{geometry}
\usepackage{amsmath,amssymb,mathtools}
\usepackage{hyperref}
\usepackage{titlesec}
\usepackage{enumitem}

% ---------- Macros ----------
\newcommand{\ee}{\mathrm{e}}
\newcommand{\dd}{\mathrm{d}}

% ---------- Compact spacing ----------
\setlength{\parindent}{0pt}
\setlength{\parskip}{0.25em}
\setlength{\abovedisplayskip}{0.65em}
\setlength{\belowdisplayskip}{0.65em}
\setlength{\abovedisplayshortskip}{0.45em}
\setlength{\belowdisplayshortskip}{0.45em}
\setlength{\jot}{0.35em}

\titlespacing*{\section}{0pt}{1.2em}{0.45em}
\titlespacing*{\subsection}{0pt}{0.9em}{0.35em}

\setlist{nosep,leftmargin=*}

\begin{document}
	
	% ============================================================
	% Title
	% ============================================================
	
	\begin{center}
		
		{\Large\textbf{Quantum Field Theory}}
		
		\vspace{0.4cm}
		
		{\Large\textbf{First Exercise Series}}
		
		\vspace{0.7cm}
		
		{\large\textbf{Four-Dimensional Lattice Path Propagator}}
		
		\vspace{0.25cm}
		
		{\large Complete Derivation and Continuum Limit}
		
		\vspace{0.7cm}
		
		{\large\textbf{Yasaman Mahzarnia}}
		
	\end{center}
	
	\vspace{0.8cm}
	
	% ============================================================
	% Problem
	% ============================================================
	
	\section*{Problem}
	
	Consider a four-dimensional hypercubic Euclidean lattice with lattice spacing
	$\Delta$. A path moves only between nearest-neighbour lattice sites. For a path
	containing $L$ links, assign the weight
	\[
	\ee^{-m_0 L}.
	\]
	
	The lattice propagator is
	\begin{equation}
		G_\Delta(x)
		=
		\sum_{\text{paths }0\to x}
		\ee^{-m_0 L(\text{path})}.
		\label{eq:definition}
	\end{equation}
	
	The goal is to compute this path sum and determine how $m_0$ must depend on
	$\Delta$ in order to obtain a nontrivial continuum limit.
	
	% ============================================================
	% Lattice sites and allowed steps
	% ============================================================
	
	\section{Lattice Sites and Allowed Steps}
	
	A lattice site is labelled by an integer four-vector
	\begin{equation}
		n=(n_1,n_2,n_3,n_4),
		\qquad
		n_\mu\in\mathbb{Z}.
	\end{equation}
	
	Its physical coordinate is
	\begin{equation}
		x_\mu=\Delta n_\mu,
		\qquad
		x=\Delta n.
		\label{eq:x-delta-n}
	\end{equation}
	
	The vector $n$ is dimensionless, while $x$ has dimensions of length.
	
	At every site there are two directions along each of the four axes,
	\[
	+\hat e_\mu,
	\qquad
	-\hat e_\mu,
	\qquad
	\mu=1,2,3,4.
	\]
	
	Therefore there are
	\[
	2\times4=8
	\]
	allowed nearest-neighbour steps at every site.
	
	% ============================================================
	% Grouping paths
	% ============================================================
	
	\section{Grouping Paths by Their Number of Links}
	
	Let $C_L(n)$ be the number of paths with exactly $L$ steps from $0$ to $n$.
	
	All such paths have the same weight $\ee^{-m_0L}$, hence
	\begin{equation}
		\boxed{
			G_\Delta(n)
			=
			\sum_{L=0}^{\infty}
			\ee^{-m_0L}C_L(n).
		}
		\label{eq:G-C}
	\end{equation}
	
	The initial condition is
	\begin{equation}
		\boxed{
			C_0(n)=\delta_{n,0}.
		}
		\label{eq:initial-condition}
	\end{equation}
	
	With zero steps, the walker is at the origin only.
	
	% ============================================================
	% Long paths and critical value
	% ============================================================
	
	\section{Long Paths and the Critical Value}
	
	An individual path of length $L$ has weight $\ee^{-m_0L}$.
	
	However, if the endpoint is unrestricted, there are $8^L$ possible step
	sequences. This does not mean that $8^L$ paths arrive at a fixed endpoint;
	for fixed $n$,
	\[
	C_L(n)\leq 8^L.
	\]
	
	Summing over all endpoints gives
	\begin{equation}
		\sum_n C_L(n)=8^L.
	\end{equation}
	
	Thus the total weight of all unrestricted paths of length $L$ scales as
	\begin{equation}
		W_L
		\sim
		8^L\ee^{-m_0L}
		=
		\left(8\ee^{-m_0}\right)^L
		=
		\ee^{-(m_0-\ln8)L}.
		\label{eq:entropy-energy}
	\end{equation}
	
	Therefore,
	\begin{align}
		m_0>\ln8
		&\quad\Longrightarrow\quad
		\text{long paths are exponentially suppressed},
		\\
		m_0=\ln8
		&\quad\Longrightarrow\quad
		\text{critical boundary of convergence},
		\\
		m_0<\ln8
		&\quad\Longrightarrow\quad
		\text{the path sum diverges}.
	\end{align}
	
	For the continuum limit, the critical value is approached from above:
	\begin{equation}
		m_0(\Delta)\to\ln8^+.
	\end{equation}
	
	% ============================================================
	% Recursion
	% ============================================================
	
	\section{Recursion Relation for Path Counting}
	
	To arrive at $n$ after $L+1$ steps, the final step must come from one of the
	eight nearest neighbours of $n$.
	
	Hence
	\begin{equation}
		\boxed{
			C_{L+1}(n)
			=
			\sum_{\mu=1}^{4}
			\left[
			C_L(n-\hat e_\mu)
			+
			C_L(n+\hat e_\mu)
			\right].
		}
		\label{eq:recurrence}
	\end{equation}
	
	% ============================================================
	% Fourier transform
	% ============================================================
	
	\section{Fourier Transform on the Lattice}
	
	We Fourier transform with respect to the lattice coordinate $n$, rather than
	the path length $L$, because a lattice shift becomes multiplication by a phase
	in Fourier space.
	
	Define
	\begin{equation}
		\boxed{
			\widetilde C_L(k)
			=
			\sum_{n\in\mathbb{Z}^4}
			\ee^{-ik\cdot n}C_L(n).
		}
		\label{eq:fourier-forward}
	\end{equation}
	
	The inverse transform is
	\begin{equation}
		\boxed{
			C_L(n)
			=
			\int_{[-\pi,\pi]^4}
			\frac{\dd^4k}{(2\pi)^4}
			\,
			\ee^{ik\cdot n}
			\widetilde C_L(k).
		}
		\label{eq:fourier-inverse}
	\end{equation}
	
	Here
	\[
	k\cdot n
	=
	\sum_{\mu=1}^{4}k_\mu n_\mu.
	\]
	
	Since $n_\mu$ is an integer,
	\[
	\ee^{i(k_\mu+2\pi)n_\mu}
	=
	\ee^{ik_\mu n_\mu},
	\]
	so it is sufficient to take
	\[
	-\pi\leq k_\mu\leq\pi.
	\]
	
	% ============================================================
	% Initial condition in Fourier space
	% ============================================================
	
	\section{Initial Condition in Fourier Space}
	
	At $L=0$, insert $C_0(n)=\delta_{n,0}$ into the Fourier-transform definition:
	\begin{align}
		\widetilde C_0(k)
		&=
		\sum_n
		\ee^{-ik\cdot n}C_0(n)
		\\
		&=
		\sum_n
		\ee^{-ik\cdot n}\delta_{n,0}
		\\
		&=
		\ee^{-ik\cdot0}
		=
		1.
	\end{align}
	
	Thus
	\begin{equation}
		\boxed{
			\widetilde C_0(k)=1.
		}
		\label{eq:C0tilde}
	\end{equation}
	
	The Kronecker delta selects only the origin term.
	
	% ============================================================
	% Solving recursion
	% ============================================================
	
	\section{Solving the Recursion Relation in Fourier Space}
	
	Fourier transforming Eq.~\eqref{eq:recurrence}, use the changes of variables
	$r=n-\hat e_\mu$ and $r=n+\hat e_\mu$, respectively:
	\begin{align}
		\sum_n
		\ee^{-ik\cdot n}C_L(n-\hat e_\mu)
		&=
		\ee^{-ik_\mu}\widetilde C_L(k),
		\\
		\sum_n
		\ee^{-ik\cdot n}C_L(n+\hat e_\mu)
		&=
		\ee^{ik_\mu}\widetilde C_L(k).
	\end{align}
	
	Therefore,
	\begin{align}
		\widetilde C_{L+1}(k)
		&=
		\sum_{\mu=1}^{4}
		\left(
		\ee^{-ik_\mu}
		+
		\ee^{ik_\mu}
		\right)
		\widetilde C_L(k)
		\\
		&=
		\left[
		2\sum_{\mu=1}^{4}\cos k_\mu
		\right]
		\widetilde C_L(k).
		\label{eq:recurrence-fourier}
	\end{align}
	
	We used
	\[
	\ee^{i\theta}+\ee^{-i\theta}=2\cos\theta.
	\]
	
	The factor
	\[
	2\sum_\mu\cos k_\mu
	\]
	is the one-step hopping factor in Fourier space.
	
	% ============================================================
	% Solution for C_L
	% ============================================================
	
	\section{Solution for $C_L(n)$}
	
	Define
	\begin{equation}
		F(k)
		=
		2\sum_{\mu=1}^{4}\cos k_\mu.
	\end{equation}
	
	Since $\widetilde C_0(k)=1$, Eq.~\eqref{eq:recurrence-fourier} gives
	\begin{equation}
		\boxed{
			\widetilde C_L(k)
			=
			F(k)^L
			=
			\left[
			2\sum_{\mu=1}^{4}\cos k_\mu
			\right]^L.
		}
		\label{eq:CLtilde-solution}
	\end{equation}
	
	Applying the inverse transform,
	\begin{equation}
		\boxed{
			C_L(n)
			=
			\int_{[-\pi,\pi]^4}
			\frac{\dd^4k}{(2\pi)^4}
			\,
			\ee^{ik\cdot n}
			\left[
			2\sum_{\mu=1}^{4}\cos k_\mu
			\right]^L.
		}
		\label{eq:CL-solution}
	\end{equation}
	
	% ============================================================
	% Exact lattice propagator
	% ============================================================
	
	\section{Exact Lattice Propagator}
	
	Insert Eq.~\eqref{eq:CL-solution} into Eq.~\eqref{eq:G-C}:
	\begin{align}
		G_\Delta(n)
		&=
		\int_{[-\pi,\pi]^4}
		\frac{\dd^4k}{(2\pi)^4}
		\,
		\ee^{ik\cdot n}
		\sum_{L=0}^{\infty}
		\left[
		2\ee^{-m_0}
		\sum_{\mu=1}^{4}\cos k_\mu
		\right]^L.
	\end{align}
	
	The sum over $L$ is geometric:
	\begin{equation}
		\sum_{L=0}^{\infty}r^L
		=
		\frac{1}{1-r},
		\qquad
		|r|<1.
	\end{equation}
	
	Hence
	\begin{equation}
		\boxed{
			G_\Delta(n)
			=
			\int_{[-\pi,\pi]^4}
			\frac{\dd^4k}{(2\pi)^4}
			\,
			\frac{
				\ee^{ik\cdot n}
			}{
				1
				-
				2\ee^{-m_0}
				\displaystyle\sum_{\mu=1}^{4}\cos k_\mu
			}.
		}
		\label{eq:lattice-propagator}
	\end{equation}
	
	% ============================================================
	% Convergence
	% ============================================================
	
	\section{Convergence Condition}
	
	The largest value of
	\[
	2\sum_\mu\cos k_\mu
	\]
	is at $k=0$:
	\begin{equation}
		2\sum_{\mu=1}^{4}\cos0=8.
	\end{equation}
	
	Thus convergence requires
	\begin{equation}
		8\ee^{-m_0}<1,
	\end{equation}
	or
	\begin{equation}
		\boxed{
			m_0>\ln8.
		}
		\label{eq:convergence}
	\end{equation}
	
	% ============================================================
	% Continuum variables
	% ============================================================
	
	\section{Continuum Variables}
	
	Take $\Delta\to0$ while keeping $x=\Delta n$ fixed.
	
	The lattice plane wave is $\ee^{ik\cdot n}$, while the continuum plane wave is
	$\ee^{ip\cdot x}$.
	
	Matching their phases gives
	\begin{equation}
		k\cdot n
		=
		p\cdot x
		=
		p\cdot(\Delta n),
	\end{equation}
	and hence
	\begin{equation}
		\boxed{
			k=\Delta p.
		}
		\label{eq:k-delta-p}
	\end{equation}
	
	Therefore
	\begin{equation}
		\ee^{ik\cdot n}
		=
		\ee^{ip\cdot x},
		\qquad
		\dd^4k
		=
		\Delta^4\dd^4p.
	\end{equation}
	
	Equation~\eqref{eq:lattice-propagator} becomes
	\begin{equation}
		G_\Delta(x)
		=
		\int_{[-\pi/\Delta,\pi/\Delta]^4}
		\frac{\Delta^4\dd^4p}{(2\pi)^4}
		\,
		\frac{
			\ee^{ip\cdot x}
		}{
			1
			-
			2\ee^{-m_0}
			\displaystyle\sum_{\mu=1}^{4}
			\cos(\Delta p_\mu)
		}.
		\label{eq:propagator-p}
	\end{equation}
	
	% ============================================================
	% Small Delta expansion
	% ============================================================
	
	\section{Small-$\Delta$ Expansion}
	
	For fixed physical momentum $p$,
	\begin{equation}
		\cos(\Delta p_\mu)
		=
		1
		-
		\frac{\Delta^2p_\mu^2}{2}
		+
		O(\Delta^4p_\mu^4).
	\end{equation}
	
	With
	\[
	p^2
	=
	\sum_{\mu=1}^{4}p_\mu^2,
	\]
	this gives
	\begin{equation}
		2\sum_{\mu=1}^{4}\cos(\Delta p_\mu)
		=
		8
		-
		\Delta^2p^2
		+
		O(\Delta^4p^4).
	\end{equation}
	
	Hence the denominator is
	\begin{align}
		D_\Delta(p)
		&=
		1
		-
		2\ee^{-m_0}
		\sum_{\mu=1}^{4}\cos(\Delta p_\mu)
		\\
		&=
		1
		-
		8\ee^{-m_0}
		+
		\ee^{-m_0}\Delta^2p^2
		+
		O(\Delta^4p^4).
		\label{eq:denominator-expansion}
	\end{align}
	
	% ============================================================
	% Tuning m0
	% ============================================================
	
	\section{Why $m_0$ Must Be Tuned}
	
	If $m_0>\ln8$ is held fixed, then
	\[
	1-8\ee^{-m_0}
	\]
	is a nonzero constant, while the $p^2$ term vanishes as $\Delta^2$.
	
	The continuum limit would then lose the momentum dependence needed for a
	propagator of the form
	\[
	\frac{1}{p^2+m^2}.
	\]
	
	Therefore,
	\begin{equation}
		\boxed{
			1-8\ee^{-m_0(\Delta)}
			=
			O(\Delta^2).
		}
	\end{equation}
	
	This implies
	\[
	m_0(\Delta)\to\ln8^+.
	\]
	
	Define the physical mass $m$ by choosing
	\begin{equation}
		\boxed{
			1-8\ee^{-m_0(\Delta)}
			=
			\ee^{-m_0(\Delta)}m^2\Delta^2.
		}
		\label{eq:tuning-condition}
	\end{equation}
	
	Multiplying by $\ee^{m_0(\Delta)}$ gives
	\begin{equation}
		\ee^{m_0(\Delta)}
		=
		8+m^2\Delta^2.
	\end{equation}
	
	Thus
	\begin{equation}
		\boxed{
			m_0(\Delta)
			=
			\ln(8+m^2\Delta^2)
			=
			\ln8
			+
			\frac{m^2\Delta^2}{8}
			+
			O(\Delta^4).
		}
		\label{eq:m0-tuning}
	\end{equation}
	
	% ============================================================
	% Continuum propagator
	% ============================================================
	
	\section{Continuum Klein--Gordon Propagator}
	
	With this tuning,
	\begin{equation}
		D_\Delta(p)
		=
		\ee^{-m_0(\Delta)}
		\Delta^2(p^2+m^2)
		+
		O(\Delta^4p^4).
	\end{equation}
	
	Therefore,
	\begin{align}
		G_\Delta(x)
		&\simeq
		\ee^{m_0(\Delta)}\Delta^2
		\int
		\frac{\dd^4p}{(2\pi)^4}
		\,
		\frac{\ee^{ip\cdot x}}{p^2+m^2}.
		\label{eq:G-asymptotic}
	\end{align}
	
	The raw path sum contains the overall lattice-dependent factor
	$\ee^{m_0}\Delta^2$.
	
	Remove it by defining
	\begin{equation}
		\boxed{
			D_E(x)
			=
			\lim_{\Delta\to0}
			\frac{\ee^{-m_0(\Delta)}}{\Delta^2}
			G_\Delta(x).
		}
		\label{eq:normalization}
	\end{equation}
	
	Using Eq.~\eqref{eq:G-asymptotic},
	\begin{equation}
		\boxed{
			D_E(x)
			=
			\int_{\mathbb{R}^4}
			\frac{\dd^4p}{(2\pi)^4}
			\,
			\frac{\ee^{ip\cdot x}}{p^2+m^2}.
		}
		\label{eq:KG-propagator}
	\end{equation}
	
	This is the Euclidean propagator of a free Klein--Gordon field.
	
	% ============================================================
	% Final result
	% ============================================================
	
	\section*{Final Result}
	
	The exact lattice propagator is
	\begin{equation}
		\boxed{
			G_\Delta(n)
			=
			\int_{[-\pi,\pi]^4}
			\frac{\dd^4k}{(2\pi)^4}
			\,
			\frac{
				\ee^{ik\cdot n}
			}{
				1
				-
				2\ee^{-m_0}
				\displaystyle\sum_{\mu=1}^{4}\cos k_\mu
			}.
		}
	\end{equation}
	
	The necessary tuning is
	\begin{equation}
		\boxed{
			m_0(\Delta)
			=
			\ln(8+m^2\Delta^2)
			=
			\ln8
			+
			\frac{m^2\Delta^2}{8}
			+
			O(\Delta^4).
		}
	\end{equation}
	
	Finally,
	\begin{equation}
		\boxed{
			D_E(x)
			=
			\lim_{\Delta\to0}
			\frac{\ee^{-m_0(\Delta)}}{\Delta^2}
			G_\Delta(x)
			=
			\int_{\mathbb{R}^4}
			\frac{\dd^4p}{(2\pi)^4}
			\,
			\frac{\ee^{ip\cdot x}}{p^2+m^2}.
		}
	\end{equation}

\newpage
\subsection*{AI Use Declaration}

\textbf{Model(s) used:}
Perplexity AI (for generating the full LaTeX code and helping to explain the solution steps).

\textbf{Revisions/corrections made to model output, and clarifications added:}

I did not change the mathematical steps or the final formulas, because the derivation produced by the AI looked correct to me. My follow-up questions were about conceptual points that I did not fully understand at first: why we Fourier transform with respect to the lattice coordinate $n$ instead of the path length $L$, why a momentum variable $p$ appears when we set $k = \Delta p$ and try to match a continuum plane wave $e^{ip\cdot x}$, and how the tuning of $m_0(\Delta)$ is chosen so that the lattice result reproduces the known form of the continuum Klein--Gordon propagator.

The AI explained that the introduction of $p$ is motivated by the fact that we already know what the Klein--Gordon propagator looks like in momentum space and are effectively designing the lattice limit to match it. Based on these answers, I only refined the textual explanations around these points; the equations and their order were left unchanged.

\textbf{Comparisons/checks against the lecture notes:}

I did not formally compare the output against the lecture notes, but since this material was covered in the previous lecture, I checked the main steps against what I remembered from class.

\textbf{Calculation steps, limiting cases, or consistency checks verified by hand:}

I reviewed and checked almost all main calculations on paper.

\textbf{Calculation steps, limiting cases, or consistency checks subjected to automated verification (via Mathematica, Python, LLMs):}

I did not check any calculation with Mathematica or Python; I only used a text-based LLM to generate and organize the LaTeX code and explanations.


\end{document}
```
